% 当前正文配置：依据 TP-0.2；修改需同步受影响的接口与图表。
\newcommand{\RequirementHardwareBasis}{国产 RISC-V 处理器及嵌入式芯片集成}
\newcommand{\RequirementHardwareResponse}{先完成现有 FPGA 上的开源架构研究原型，国产芯片落地单独推进}
\newcommand{\RequirementHardwareVerification}{硬件来源、开源许可、联合 RTL 与板级资源验证}
\newcommand{\RequirementSoftwareBasis}{国产实时操作系统、可裁剪 Linux 及混合部署}
\newcommand{\RequirementSoftwareResponse}{RT-Thread 与 Linux 两条独立启动路径，共用推理核心源码}
\newcommand{\RequirementSoftwareVerification}{两种系统分别启动并执行同一模型}
\newcommand{\RequirementFunctionsBasis}{智能感知、机器视觉、模型推理及移动数据中心接入}
\newcommand{\RequirementFunctionsResponse}{基础推理先行，采集、通信和中心接入通过适配层扩展}
\newcommand{\RequirementFunctionsVerification}{模型结果对照，后续逐项开展实物采集与联调}
\newcommand{\RequirementDeliveryBasis}{芯片集成、驱动、移植与感控应用的完整技术体系}
\newcommand{\RequirementDeliveryResponse}{交付 SoC 配置、BSP、算子库、转换器、运行时及验证说明}
\newcommand{\RequirementDeliveryVerification}{按配置、接口、应用与证据逐层复现}
